\scanpage{33}
% archive sheet number: 32
\mspage{\handwritten{31}}
% L1–3
En effet, on peut supposer $k$ algébriquement clos, de sorte que $T$ est contenu dans un tore maximal \deleted{\unsure{de $G$}} $T'$. Soit $C'$ le centralisateur de ce dernier, alors on a $C'\handwritten{\subset}C$, donc $\dim C'\handwritten{\leq}\dim C$, cqfd.

% L4–8
b) \deleted{\unsure{Le raisonnement de a) nous montre que}} Si $\handwritten{\rho}_{r}$ est locqlement constant, alors pour tout tore $T$ dans $G$, \deleted{\unsure{et}} et tout $s\handwritten{\in}S$, si $T_{s}$ est \inserted{tore} maximal dans $G_{s}$, alors il existe un voisinage ouvert $U$ de $s$ tel que $\deleted{\unsure{G}}T\handwritten{\vert}U$ soit tore maximal dans $G\ U$. Utilisant maintenant le raisonnement de a), on voit que \deleted{\unsure{si}} (i)$\handwritten{\Rightarrow}$(ii bis). D'autre part (ii bis)$\handwritten{\Rightarrow}$(i),
% L9–15
car \deleted{on peut supposer $S$} lorsque $G$ admet un tore maximal $T$, alors il est évident que $\handwritten{\rho}_{r}(s)=\dim T_{s}$ est une fonction localement constante de \deleted{\unsure{$s$}} $s$, or nous avons signalé au début de la démonstration de a) que la question de la continuité de $\handwritten{\rho}_{r}$ était locale pour la topologie fidèlement plate. D\deleted{\unsure{v}}onc (i)$\handwritten{\Leftrightarrow}$(ii bis), et évidemment (ii)$\handwritten{\Rightarrow}$(ii bis), reste à prouver l'implication inverse \inserted{(i)$\handwritten{\Longrightarrow}$(ii)}.
% L15–19
Pour ceci, introduisons le foncteur $F$ de \inserted{\handwritten{Exp \uline{XI}}}4.1. , qui est donc un préschéma \deleted{\unsure{localement de}} lisse \inserted{\handwritten{et séparé}} sur $S$, et considérons le sous-foncteur $\deleted{\unsure{F}}\handwritten{\mathscr{T}}$ de $F$, dont la valeur \deleted{sur} pour un $\handwritten{S'}$ sur $S$ est l'ensemble des tores maximaux dans $G_{\handwritten{S'}}$. Utilisant la remarque précédente \deleted{\unsure{que}} \inserted{que\handwritten{,} moyennant (i)\handwritten{,}} un tore dans ${G\deleted{\unsure{s|a}}}_{\handwritten{S'}}$ qui est maximal dans la fibre d'un
% L20–24
point $\handwritten{s\in S'}$ est maximal au dessus d\deleted{\unsure{e}}'un voisinage ouvert de $\handwritten{s}$, on voit que $\handwritten{\mathscr{T}}$ \margin{\handwritten{$\mathscr{T}$}}est représentable par un sous-préschéma ouvert de $F$, et par suite est lisse \inserted{\handwritten{et séparé}} sur $S$. Comme le morphisme structural $\handwritten{\mathscr{T}\to S}$ est évidem\deleted{\unsure{-}}ment surjectif, il admet donc \deleted{localement} une section localement au sens de la topologie étale en vertu de \inserted{\handwritten{Exp \uline{XI}}}1.1\handwritten{0}. , ce qui \handwritten{p}rouve (i)$\handwritten{\Rightarrow}$(ii).

% L25–26
c) C'est une conséquence immédiate de \inserted{\handwritten{Exp \uline{XI}}}5.4. bis, compte tenu du théorème de conjugaison de Borel rappelé au début du numéro.

% L27–28
d) Moyennant la remarque du début de la démonstration dans a), et compte tenu de b), on peut supposer qu'il existe un tore maximal $T$
